\sheet{9}{Z-Transform}{Chapter~9 ($z$-transform, transfer functions of discrete-time systems)}

\begin{goals}
\begin{itemize}
  \item Compute $z$-transforms with their region of convergence (ROC) and
    invert rational $X(z)$ by partial fractions for causal, anticausal and
    two-sided sequences.
  \item Derive $H(z)$, poles and zeros of the RPM smoothers of Sheet~3 and decide
    stability from the pole positions.
  \item Read the magnitude response from the pole-zero map (geometric
    interpretation) and design a two-pole resonator for the knock frequency.
  \item Build a pole-zero explorer (polynomial roots, ASCII pole-zero map,
    frequency and impulse response).
  \item Implement the resonator in Q15 fixed point on the Arduino Uno and
    observe coefficient quantisation, zero-input limit cycles and overflow
    oscillations for $r\to1$.
\end{itemize}
\end{goals}

\section*{Background}

The (bilateral) $z$-transform of $x[n]$ is $X(z)=\sum_{n=-\infty}^{\infty}x[n]z^{-n}$,
defined in a ring-shaped ROC $r_1<|z|<r_2$. An LTI system with the difference
equation $\sum_k a_ky[n-k]=\sum_k b_kx[n-k]$ has the transfer function
\[
  H(z)=\frac{b_0+b_1z^{-1}+\dots+b_Mz^{-M}}{a_0+a_1z^{-1}+\dots+a_Nz^{-N}}
      =\frac{b_0}{a_0}\,z^{N-M}\,\frac{\prod_k(z-z_k)}{\prod_k(z-p_k)} .
\]
A causal system is stable iff all poles lie inside the unit circle (then the
ROC $|z|>\max|p_k|$ contains the unit circle), and its frequency response is
$H(\e^{\jj\Omega})$, $\Omega=2\pi f/f_s$. Geometrically,
$|H(\e^{\jj\Omega})|=\bigl|\frac{b_0}{a_0}\bigr|\prod_k|\e^{\jj\Omega}-z_k|\big/\prod_k|\e^{\jj\Omega}-p_k|$:
a pole close to the unit circle produces a peak, a zero on it a notch.

The \textbf{two-pole resonator} with poles $p_{1,2}=r\e^{\pm\jj\theta}$,
$\theta=2\pi f_0/f_s$,
\begin{equation}
  H(z)=\frac{b_0}{1+a_1z^{-1}+a_2z^{-2}},\qquad a_1=-2r\cos\theta,\quad a_2=r^2,
  \label{eq:s9res}
\end{equation}
is the simplest digital knock filter. Throughout this sheet
$f_0=\SI{6.5}{kHz}$ and $f_s=\SI{40}{kHz}$ ($\theta=58.5^\circ$).

%-------------------------------------------------------------------------
\begin{exercise}{Analytical $z$-transforms}{\Paper}
\begin{tasks}
  \item Compute the $z$-transform and the ROC of $\delta[n-k]$, $\varepsilon[n]$,
    $a^n\varepsilon[n]$, $-a^n\varepsilon[-n-1]$ and
    $r^n\cos(\theta n)\,\varepsilon[n]$. Why is the ROC indispensable?
  \item Find all sequences with
    $X(z)=\dfrac{1}{(1-0.5z^{-1})(1-0.8z^{-1})}$ by partial fractions (one for
    each possible ROC). Which one is causal, which one has a DTFT? Check
    $x[0]$ and $x[1]$ of the causal one with the difference equation.
  \item \textbf{RPM smoothers of Sheet~3.} Derive $H(z)$, poles and zeros of
    \begin{tasks}
      \item the exponential smoother $y[n]=a\,y[n-1]+(1-a)\,x[n]$,
      \item the $M$-point moving average.
    \end{tasks}
    When are they stable? Give $|H|$ at $\Omega=0$ and $\Omega=\pi$ for $a=0.9$
    and the notch frequencies of the moving average for $M=8$ at
    $f_s=\SI{40}{kHz}$. What happens to the pole at $z=1$ of
    $(1-z^{-M})/(1-z^{-1})$?
  \item \textbf{Resonator.} Show that \eqref{eq:s9res} has the poles
    $r\e^{\pm\jj\theta}$. Derive the stability triangle in the
    $(a_1,a_2)$-plane ($|a_2|<1$, $|a_1|<1+a_2$) and mark the resonator.
  \item Using the geometric interpretation, show that the $-3$\,dB bandwidth is
    $\Delta\Omega\approx2(1-r)$, i.e.\ $B\approx(1-r)f_s/\pi$, for $r$ close to 1.
    Choose $r$ for $B=\SI{400}{Hz}$. Compare with the impulse-invariant image
    $r=\e^{-T/\tau}$ of the analog knock resonator of Sheet~7
    ($\tau=\SI{0.8}{ms}$).
  \item Determine $b_0$ such that $|H(\e^{\jj\theta})|=1$ and evaluate it for
    $r=0.97$.
\end{tasks}
\end{exercise}

\begin{solution}
a) $\delta[n-k]\rightarrow z^{-k}$ (all $z$, except $z=0$ for $k>0$ and
$z=\infty$ for $k<0$). $\varepsilon[n]\rightarrow\frac{1}{1-z^{-1}}$, $|z|>1$.
$a^n\varepsilon[n]\rightarrow\sum_{n\ge0}(az^{-1})^n=\frac{1}{1-az^{-1}}$, $|z|>|a|$.
$-a^n\varepsilon[-n-1]\rightarrow-\sum_{m\ge1}(a^{-1}z)^m=\frac{1}{1-az^{-1}}$,
$|z|<|a|$ -- the same algebraic expression: only the ROC tells which sequence is
meant. $r^n\cos(\theta n)\varepsilon[n]=\frac{r^n}2(\e^{\jj\theta n}+\e^{-\jj\theta n})\varepsilon[n]
\rightarrow\frac{1-r\cos\theta\,z^{-1}}{1-2r\cos\theta\,z^{-1}+r^2z^{-2}}$, $|z|>r$.

b) $X(z)=\frac{A}{1-0.5z^{-1}}+\frac{B}{1-0.8z^{-1}}$ with
$A=\frac{1}{1-0.8/0.5}=-\frac53$, $B=\frac{1}{1-0.5/0.8}=\frac83$.
\begin{itemize}
  \item $|z|>0.8$: $x[n]=\bigl(-\frac53\,0.5^n+\frac83\,0.8^n\bigr)\varepsilon[n]$
    (causal, ROC contains the unit circle $\Rightarrow$ DTFT exists).
  \item $0.5<|z|<0.8$: $x[n]=-\frac53\,0.5^n\varepsilon[n]-\frac83\,0.8^n\varepsilon[-n-1]$
    (two-sided, no DTFT).
  \item $|z|<0.5$: $x[n]=\frac53\,0.5^n\varepsilon[-n-1]-\frac83\,0.8^n\varepsilon[-n-1]$
    (anticausal, no DTFT).
\end{itemize}
Check: $y[n]=1.3y[n-1]-0.4y[n-2]+\delta[n]$ gives $y[0]=1$, $y[1]=1.3$;
the formula: $-\frac53+\frac83=1$, $-\frac56+\frac{32}{15}=1.3$.

c) i.~$H(z)=\frac{1-a}{1-az^{-1}}=\frac{(1-a)z}{z-a}$: zero at 0, pole at $a$;
stable iff $|a|<1$. $|H(\e^{\jj0})|=1$, $|H(\e^{\jj\pi})|=\frac{1-a}{1+a}=0.0526$
($\SI{-25.6}{dB}$ for $a=0.9$).
ii.~$H(z)=\frac1M\sum_{k=0}^{M-1}z^{-k}=\frac1M\frac{1-z^{-M}}{1-z^{-1}}=\frac{z^M-1}{M\,z^{M-1}(z-1)}$.
The roots of $z^M=1$ are $\e^{\jj2\pi k/M}$; the one at $z=1$ cancels the pole at
$z=1$, so there remain $M-1$ zeros on the unit circle and $M-1$ poles at
$z=0$: an FIR filter, always stable. For $M=8$: notches at
$kf_s/8=5,10,15,\SI{20}{kHz}$.

d) $z^2+a_1z+a_2=z^2-2r\cos\theta\,z+r^2=(z-r\e^{\jj\theta})(z-r\e^{-\jj\theta})$.
Stability triangle: for complex poles $|p|^2=a_2<1$; for real poles the
polynomial $P(z)=z^2+a_1z+a_2$ must have both roots in $(-1,1)$, i.e.\
$P(1)=1+a_1+a_2>0$ and $P(-1)=1-a_1+a_2>0$ $\Rightarrow|a_1|<1+a_2$, $|a_2|<1$.
The resonator with $r=0.97$ ($a_1=-1.0136$, $a_2=0.9409$) lies inside, close
to the upper edge $a_2=1$.

e) For $\Omega$ near $\theta$ the distance to the far pole is almost constant,
$|\e^{\jj\theta}-r\e^{-\jj\theta}|\approx2\sin\theta$; the distance to the near
pole is $|\e^{\jj\Omega}-r\e^{\jj\theta}|\approx\sqrt{(1-r)^2+(\Omega-\theta)^2}$.
It grows by $\sqrt2$ for $|\Omega-\theta|=1-r$: $\Delta\Omega=2(1-r)$,
$B=\frac{\Delta\Omega}{2\pi}f_s=(1-r)f_s/\pi$. $B=\SI{400}{Hz}$:
$r=1-\pi B/f_s=0.9686$. Impulse invariance: $r=\e^{-T/\tau}=\e^{-0.03125}=0.9692$,
$B=\SI{392}{Hz}$, close to the analog \SI{398}{Hz} -- both views agree, because
$1-\e^{-T/\tau}\approx T/\tau$ and $B_\text{analog}=1/(\pi\tau)$.

f) $|H(\e^{\jj\theta})|=\frac{b_0}{|1-r\e^{\jj\theta}\e^{-\jj\theta}|\,|1-r\e^{-\jj\theta}\e^{-\jj\theta}|}
=\frac{b_0}{(1-r)\,|1-r\e^{-2\jj\theta}|}$, hence
$b_0=(1-r)\sqrt{1-2r\cos2\theta+r^2}$; $r=0.97$: $b_0=0.050393$. Peak gain
without normalisation ($b_0=1$): $1/0.0504=19.8$ (\SI{26}{dB}) -- important for
the fixed-point scaling in Exercise~9.3.
\end{solution}

%-------------------------------------------------------------------------
\begin{exercise}{Pole-zero explorer}{\JSLinux}
Program \codefile{jslinux/sheet09/pz\_explorer.c} reads the coefficients from
the command line,
\begin{lstlisting}[style=shell]
./pz_explorer -fs 40000 -b 0.1 -a 1 -0.9      # exponential smoother
./pz_explorer -fs 40000 -res 6500 0.97        # resonator, gain 1 at f0
\end{lstlisting}
computes poles and zeros, prints them in polar form (radius, angle in Hz),
decides stability, draws the pole-zero map with \texttt{ap\_pz()} and plots
magnitude (dB) and phase over $f$ and the impulse response.
\begin{tasks}
  \item Poles and zeros are the roots of $z^LB(z)$ and $z^LA(z)$,
    $L=\max(M,N)$ (\texttt{tf\_roots()} is given). Implement the
    Durand--Kerner iteration \texttt{poly\_roots()} for a polynomial of
    degree $n$ with $c_0=1$:
    \[ r_i\leftarrow r_i-\frac{c(r_i)}{\prod_{j\ne i}(r_i-r_j)},\qquad
       r_i^{(0)}=(0.4+0.9\jj)^{i+1}. \]
    Use the small complex struct of the template (no \texttt{complex.h}).
  \item Implement the geometric evaluation \texttt{freq\_geom()} and the
    impulse response from the difference equation. The program compares the
    geometric magnitude with the direct evaluation of $B/A$.
  \item Explore and report (pole-zero map, peak/notch, bandwidth, $|h[n]|$):
    \begin{tasks}
      \item exponential smoother $a=0.9$; moving average $M=8$;
      \item resonator $r=0.9$, $0.97$, $0.99$: compare the $-3$\,dB bandwidth
        with $(1-r)f_s/\pi$;
      \item resonator $r=1$ and $r=1.02$: what does the program say, what does the
        impulse response do?
      \item a notch filter at \SI{6.5}{kHz} with zeros at $\e^{\pm\jj\theta}$
        and poles at $0.95\e^{\pm\jj\theta}$. Give the coefficients and the
        notch width; why is $|H|\approx1$ far from the notch?
    \end{tasks}
  \item Why may the magnitude plot of an unstable system be misleading?
\end{tasks}
\end{exercise}

\begin{solution}
a)/b) Key parts of \codefile{solutions/jslinux/sheet09/pz\_explorer.c}:
\inputcode[firstline=53,lastline=95]{solutions/jslinux/sheet09/pz_explorer.c}
\inputcode[firstline=109,lastline=117]{solutions/jslinux/sheet09/pz_explorer.c}
The geometric and the direct magnitude agree to $10^{-14}$\,dB (notch:
$10^{-8}$\,dB, relative error near the zero). Durand--Kerner converges
quadratically for simple roots; for multiple roots (e.g.\ the $M-1$ poles at
$z=0$ of the moving average, here found exactly by \texttt{tf\_roots}) it is
only linear.

c) Program results ($f_s=\SI{40}{kHz}$, grid \SI{10}{Hz}):
\begin{itemize}
  \item Smoother $a=0.9$: pole $0.9$, zero $0$; $|H|$ \SI{0}{dB} at DC,
    \SI{-25.58}{dB} at $f_s/2$; $h[n]=0.1\cdot0.9^n$ ($|h[59]|=2.0\cdot10^{-4}$).
    Moving average $M=8$: zeros at $\pm5$, $\pm10$, $\pm15$, \SI{20}{kHz} on
    the unit circle, 7 poles at 0; $h[n]=1/8$ for $n=0..7$, then 0.
  \item Resonator ($b_0$ normalised, poles at exactly \SI{6500.0}{Hz}):
    \begin{center}\small
    \begin{tabular}{rrrrr}
    \toprule
    $r$ & $B$ measured & $(1-r)f_s/\pi$ & $|H|$ at DC & $|H|$ at $f_s/2$\\
    \midrule
    0.90 & \SI{1370}{Hz} & \SI{1273}{Hz} & \SI{-14.6}{dB} & \SI{-24.6}{dB}\\
    0.97 & \SI{400}{Hz} & \SI{382}{Hz} & \SI{-25.3}{dB} & \SI{-35.4}{dB}\\
    0.99 & \SI{140}{Hz} & \SI{127}{Hz} & \SI{-34.9}{dB} & \SI{-45.0}{dB}\\
    \bottomrule
    \end{tabular}
    \end{center}
    The approximation is good for $r\to1$ (the remaining difference for
    $r=0.99$ is the \SI{10}{Hz} grid). Smaller $B$ means a longer impulse
    response: decay $\propto r^n$, time constant $-1/(f_s\ln r)$ = \SI{0.82}{ms}
    for $r=0.97$ (cf.\ $\tau=\SI{0.8}{ms}$ of the knock burst).
  \item $r=1$: ``MARGINALLY STABLE (not BIBO stable)'', $h[n]=\sin((n+1)\theta)$,
    a sinusoidal oscillator with constant amplitude; the magnitude plot shows
    an (on the \SI{10}{Hz} grid finite) infinite peak. $r=1.02$: ``UNSTABLE'',
    $|h[n]|$ grows like $1.02^n$ ($|h[71]|=0.157$ starting from 0.034).
  \item Notch: $b=\{1,-2\cos\theta,1\}=\{1,-1.044997,1\}$,
    $a=\{1,-1.9\cos\theta,0.9025\}=\{1,-0.992748,0.9025\}$. Zeros
    $0.5225\pm0.8526\jj$ ($|z|=1$), poles $0.4964\pm0.8100\jj$ ($|p|=0.95$),
    $|H|=-116$\,dB at \SI{6.5}{kHz}, $+0.42$\,dB at DC, $+0.44$\,dB at $f_s/2$;
    $-3$\,dB stop band \SI{6200}{}--\SI{6800}{Hz}, width \SI{600}{Hz}
    $\approx(1-0.95)f_s/\pi=\SI{637}{Hz}$. Far from $\theta$ each zero and
    its neighbouring pole have almost the same distance to $\e^{\jj\Omega}$,
    the ratios cancel ($|H|\approx1$).
\end{itemize}

d) For poles outside the unit circle the ROC of the \emph{causal} system does
not contain the unit circle; $H(\e^{\jj\Omega})$ is then a formal evaluation
of the rational function, not the response to a sinusoid (which grows without
bound). The program prints a warning in this case.
\end{solution}

%-------------------------------------------------------------------------
\begin{exercise}{Digital resonator in float and Q15}{\Wokwi\ \JSLinux}
The resonator \eqref{eq:s9res} is implemented twice in the header
\texttt{resonator.h}: in \texttt{float} and in fixed point with signals in Q15
(\texttt{int16\_t}, $1\,\mathrm{LSB}=2^{-15}$), coefficients in Q14 and an
\texttt{int32\_t} accumulator. The \emph{same file} is used on JSLinux and
on the Uno; you find it in both folders\\
\codefile{jslinux/sheet09/} and \codefile{wokwi/sheet09\_resonator/}.
\begin{tasks}
  \item \textbf{(paper)} Why can $a_1$ not be stored in Q15 here? In which
    format is the accumulator, and how many bits does it need at most? Which
    quantisation step does $a_2=r^2$ get, and how large is the resulting
    change of the pole radius for $r=0.9999$?
  \item Implement the functions \texttt{res\_design}, \texttt{res\_f\_step}
    and \texttt{res\_q\_step} (rounding or truncation, saturation or
    wrap-around, and count the overflows). Note: \texttt{int} has 16 bits on
    the AVR -- cast to \texttt{int32\_t} \emph{before} multiplying. Run the
    program \texttt{limitcycle.c} and compare quantised and exact pole
    positions (part (a) of the program) and the impulse responses (b).
  \item \textbf{Zero-input limit cycles.} Starting from $y[-1]=0.5$ with
    $x=0$, the float resonator decays to zero. Part (c) of the program
    measures the amplitude of the Q15 output after \SI{5}{s}. For rounding,
    a limit cycle can persist as long as rounding maps $a_2y$ back to $y$:
    derive the dead-band bound $|y|\le\frac{0.5}{1-a_2}$\,LSB. Compare with the
    measured amplitudes for $r\to1$, rounding and truncation, and explain the
    frequency of the limit cycle.
  \item \textbf{Overflow.} Without normalisation ($b_0=1$) the gain at $f_0$ is
    about 20 (Exercise~9.1\,f). Part (d) feeds a knock burst of amplitude
    $A=0.1$ and $0.5$ into the resonator with $r=0.99$. Compare wrap-around
    and saturation. What is an \emph{overflow oscillation}?
  \item \textbf{Uno.} Open the Wokwi project (Uno, potentiometer = $r$ from
    0.90 to 1.00, switch SW1 = rounding/truncation, SW2 =
    saturation/wrap, button = next test T0/T1/T2, see the sketch header) with
    your \texttt{resonator.h}. Reproduce the results of c) and d) for
    $r\approx0.97$, $0.99$, $0.999$. Are the integer results identical to
    JSLinux? Could you also expect identical float results?
  \item Draw conclusions for a fixed-point knock filter in an ECU (choice of
    $r$, rounding, overflow handling, scaling).
\end{tasks}
\begin{hint}
\hinttext The Uno computes the test signals itself; its ADC (about
\SI{9}{kHz} maximum, Sheet~5) could not sample a \SI{6.5}{kHz} knock signal
anyway. The tests are not real-time -- we study the arithmetic; real-time
knock filtering follows on Sheet~10.
\end{hint}
\end{exercise}

\begin{solution}
a) $a_1=-2r\cos\theta=-1.0136$ for $r=0.97$ and up to $-1.045$ for $r\to1$:
outside the Q15 range $[-1,1)$. Q14 ($[-2,2)$, step $2^{-14}$) is used for all
coefficients. Products Q14$\times$Q15 are Q29; the sum of three products is
bounded by $(|b_0|+|a_1|+|a_2|)\,2^{29}<4\cdot2^{29}=2^{31}$: an
\texttt{int32\_t} is just sufficient. Step of $a_2$: $2^{-14}=6.1\cdot10^{-5}$,
$\Delta r\approx\Delta a_2/(2r)\le3.1\cdot10^{-5}$ -- comparable with
$1-r=10^{-4}$ for $r=0.9999$: the bandwidth (\SI{1.3}{Hz}) is then only
known to about $\pm30\,\%$. Program output (a):
\begin{center}\small
\begin{tabular}{rrrrrrr}
\toprule
$r$ & $b_0$ & $a_1$ & $a_2$ & Q14: $b_0,a_1,a_2$ & $r_q$ & $f_q$ [Hz]\\
\midrule
0.9 & 0.162086 & $-0.940497$ & 0.810000 & 2656, $-15409$, 13271 & 0.899999 & 6500.0\\
0.97 & 0.050393 & $-1.013647$ & 0.940900 & 826, $-16608$, 15416 & 0.970009 & 6499.9\\
0.99 & 0.016968 & $-1.034547$ & 0.980100 & 278, $-16950$, 16058 & 0.990001 & 6500.0\\
0.999 & 0.001704 & $-1.043952$ & 0.998001 & 28, $-17104$, 16351 & 0.998992 & 6500.0\\
0.9999 & 0.000171 & $-1.044892$ & 0.999800 & 3, $-17120$, 16381 & 0.999908 & 6499.9\\
\bottomrule
\end{tabular}
\end{center}
Note also the coarse $b_0$ for $r\to1$ (3 LSB: gain error up to 17\,\%).

b) Key part of \texttt{resonator.h}:
\inputcode[firstline=84,lastline=98]{solutions/wokwi/sheet09_resonator/resonator.h}
Impulse $0.5\delta[n]$, $r=0.97$, rounding: the Q15 output follows the float
output with a maximum deviation of \SI{7.2}{LSB} (peak \SI{837}{LSB}); after
400 samples the float output is $0.000$\,LSB, the Q15 output still
\SI{-2}{LSB}.

c) Zero input after $y[-1]=0.5$: a limit cycle needs $|y[n]|$ to stay constant
although the exact recursion multiplies the amplitude by $r<1$ per sample. With
rounding the oscillation behaves as if $a_2$ were 1 as long as
$|\mathrm{round}(a_2y)-a_2y|\le0.5$ can compensate the decay $(1-a_2)|y|$,
i.e.\ $(1-a_2)|y|\le0.5$ (Jackson's dead band). Measured after 200\,000 samples:
\begin{center}\small
\begin{tabular}{rrrrrrr}
\toprule
$r$ & $a_{2,q}$ & bound [LSB] & round: amp.\ [LSB] & $f$ [Hz] & trunc.: amp.\ [LSB] & $f$ [Hz]\\
\midrule
0.9 & 13271 & 2.6 & 2 & 6680 & 1 & 0 (DC)\\
0.97 & 15416 & 8.5 & 7 & 6660 & 0 & --\\
0.99 & 16058 & 25.1 & 11 & 6660 & 3 & 5720\\
0.995 & 16221 & 50.3 & 11 & 6660 & 0 & --\\
0.999 & 16351 & 248.2 & 11 & 6660 & 0 & --\\
0.9999 & 16381 & 2730.7 & 98 & 6500 & 172 & 6500\\
1.0 & 16384 & $\infty$ & 19217 & 6500 & 19207 & 6500\\
\bottomrule
\end{tabular}
\end{center}
All rounding limit cycles respect the bound; they grow with $r$ but stay far
below it for $r\ge0.99$. Their frequency is not $f_0$: the limit cycle is a
periodic integer sequence, here of period 6 samples
($f_s/6=\SI{6667}{Hz}$; the zero-crossing count gives 6660--6680). Truncation
(arithmetic shift = rounding towards $-\infty$) usually lets the output die out
but can lock into a constant $-1$\,LSB (DC limit cycle, $r=0.9$) or into
oscillations ($r=0.99$: period 7, \SI{5714}{Hz}; $r=0.9999$). Example
($r=0.99$, rounding): $\dots,-6,-11,-5,6,11,5,-6,\dots$ For $r=1$ the system is an oscillator
(amplitude $0.5/\sin\theta\cdot32768\approx19217$\,LSB), not a limit cycle.

d) $r=0.99$, $b_0=1$, burst $A\e^{-t/\tau}\sin(2\pi f_0t)$ for \SI{10}{ms}, then 0:
\begin{center}\small
\begin{tabular}{lrrr}
\toprule
case & float peak & overflows & $\max|y|$ in the last \SI{25}{ms} [LSB]\\
\midrule
$A=0.1$, wrap-around & 1.08 & 39967 & 21740\\
$A=0.1$, saturation & 1.08 & 5 & 11\\
$A=0.5$, wrap-around & 5.42 & 24 & 11\\
$A=0.5$, saturation & 5.42 & 34 & 11\\
\bottomrule
\end{tabular}
\end{center}
With two's complement wrap-around, an overflow turns a large positive value into
a large negative one; for $|a_1|+|a_2|>1$ this can sustain a full-scale
oscillation without input (\emph{overflow oscillation}): for $A=0.1$ the
filter overflows in practically every sample for the whole second. Whether it
happens depends on the trajectory ($A=0.5$ recovered by chance). Saturation
limits the error to a clipping and the filter recovers to the small rounding
limit cycle (11\,LSB).

e) The integer path uses only \texttt{int16\_t}/\texttt{int32\_t} arithmetic with
identical rounding, so the Uno must produce \emph{bit-identical} Q15 results to
JSLinux for the same $r$ (the potentiometer gives $r=0.9+0.1\cdot\mathrm{adc}/1023$,
so use the printed $r$ and Q14 coefficients for the comparison, e.g.\
$\mathrm{adc}=921\rightarrow r=0.99003$). The float path is single precision on
both, but \texttt{sinf}/\texttt{expf} of avr-libc and the evaluation order may
differ in the last bit; also a 32-bit x86 compiler (JSLinux) may keep
intermediates in 80-bit registers. Small float differences are therefore
possible, integer differences are not.

f) For $B\approx\SI{400}{Hz}$ ($r\approx0.97$) the rounding limit cycle is
7\,LSB ($\SI{-73}{dBFS}$), harmless compared with the sensor noise. Rules: use
rounding, always saturate (never wrap), scale the input with the normalised
$b_0$ (or reduce the input range) so that the output cannot exceed full scale,
store coefficients in Q14 when $|a_1|>1$, avoid $r\to1$ (narrow filters need
more bits: 32-bit state or error feedback), or use floating point on the
ESP32 (Sheet~10).
\end{solution}

\begin{deliverables}
Solutions of 9.1 (partial fractions with all ROCs, stability triangle, bandwidth
formula and $r$ for \SI{400}{Hz}); the pole-zero explorer with the screenshots
(ASCII plots) and the result table of 9.2\,c); the tables of 9.3\,a), c) and
d) from JSLinux and the corresponding Uno outputs (T0--T2 for at least two
values of $r$) with your conclusions 9.3\,f).
\end{deliverables}
